% SuperBEST_ChemBio_Catalog.tex
% SuperBEST v3 Node-Cost Catalog: Chemistry and Biology Equations
% Grand Synthesis -- Sessions Chem-1 through Bio-5
% Date: 2026-04-20
% Status: Empirical observation (computed node counts, reproducible)

\documentclass[12pt,a4paper]{article}
\usepackage{amsmath,amssymb,amsthm}
\usepackage{geometry}
\usepackage{booktabs}
\usepackage{array}
\usepackage{longtable}
\usepackage{xcolor}
\usepackage{hyperref}
\usepackage{microtype}
\usepackage{enumitem}
\geometry{margin=2.5cm}

% ── Theorem environments ────────────────────────────────────────────────────
\newtheorem{observation}{Observation}[section]
\newtheorem{theorem}{Theorem}[section]
\theoremstyle{definition}
\newtheorem{definition}[theorem]{Definition}
\newtheorem{remark}[observation]{Remark}

% ── Operator macros ─────────────────────────────────────────────────────────
\newcommand{\EML}{\operatorname{EML}}
\newcommand{\EDL}{\operatorname{EDL}}
\newcommand{\EXL}{\operatorname{EXL}}
\newcommand{\EAL}{\operatorname{EAL}}
\newcommand{\DEML}{\operatorname{DEML}}
\newcommand{\EPL}{\operatorname{EPL}}
\newcommand{\DEAL}{\operatorname{DEAL}}
\newcommand{\DEXL}{\operatorname{DEXL}}
\newcommand{\SB}{\mathrm{SB}}

% ── Column types ─────────────────────────────────────────────────────────────
\newcolumntype{L}[1]{>{\raggedright\arraybackslash}p{#1}}
\newcolumntype{C}[1]{>{\centering\arraybackslash}p{#1}}

% ════════════════════════════════════════════════════════════════════════════
\title{SuperBEST v3 Node Costs:\\
       A Catalog of Chemistry and Biology Equations\\[6pt]
       \large Grand Synthesis --- Sessions Chem-1 through Bio-5}
\author{Arturo R.\ Almaguer\\
\small monogate research \quad \texttt{almaguer1986@gmail.com}}
\date{April 2026}

\begin{document}
\maketitle

% ── Abstract ────────────────────────────────────────────────────────────────
\begin{abstract}
We present a systematic measurement of EML-family operator-node costs for
50 standard equations drawn from five areas of chemistry (rate equations,
statistical mechanics, electrochemistry, acid--base equilibria, and chemical
thermodynamics) and five areas of biology (exponential growth and decay,
population dynamics, enzyme kinetics, cooperative binding, and
spectroscopy/pharmacokinetics).  Using the SuperBEST v3 routing table
(19-node budget, 74.0\% savings over naive arithmetic), node counts range
from 2 (Malthus recursion) to 40 (MWC allosteric model, general form).
The 50 equations decompose naturally into four structural classes---pure
exponential templates, rational functions without transcendentals, log-ratio
formulas, and mixed exponential-rational forms---whose costs are determined
entirely by structural class membership rather than scientific domain.
Ten quantitative observations characterize the dominant cost patterns:
constant-folding saves up to 8 nodes per occurrence; the mechanistic--empirical
gap between MWC and Hill reaches a factor of four; linearized forms are
never cheaper than their algebraic counterparts; and the two-compartment PK
scaling law $10N-3$ generalizes across compartment number.
All results are exact and reproducible.
\end{abstract}

\tableofcontents
\newpage

% ════════════════════════════════════════════════════════════════════════════
\section{SuperBEST v3 Routing Table (Reference)}
\label{sec:routing}

The SuperBEST v3 routing table specifies the minimum number of EML-family
nodes required to compute each standard arithmetic operation.  The table
summarizes the key primitives used throughout this catalog.

\begin{table}[ht]
\centering
\caption{SuperBEST v3 --- Selected Primitive Costs}
\label{tab:routing}
\begin{tabular}{llcc}
\toprule
\textbf{Operation} & \textbf{Canonical form} & \textbf{SB nodes} & \textbf{Naive nodes} \\
\midrule
$e^x$                & \EML$(x,0)$, terminal & 1 & 1 \\
$\ln x$              & \EXL$(0,x)$, terminal & 1 & 1 \\
$x \cdot y$          & $\exp(\ln x + \ln y)$ & 3 & 7 \\
$x / y$              & $\exp(\ln x - \ln y)$ & 3 & 7 \\
$x + y$              & \EAL$(\ln x, y)$ chain & 3 & 1$^\dagger$ \\
$x - y$              & \EML$(\ln x, y)$ chain & 3 & 1$^\dagger$ \\
$x^p$ (real)         & $\exp(p \cdot \ln x)$ & 3 & 5 \\
$x^{-1}$             & \EDL$(0,x)$, terminal & 1 & 3 \\
$\exp(-x)$           & \DEML$(x,0)$, terminal & 1 & 2 \\
$x + \exp(y)$        & \EAL$(y, x)$ & 1 & 4 \\
$\ln(x) - \ln(y)$    & \EXL chain & 2 & 3 \\
$a \cdot \exp(-bx)$  & \DEML + mul & 4 & 6 \\
\bottomrule
\end{tabular}
\vspace{4pt}
\small $^\dagger$Addition is natively 1 node in standard arithmetic but costs 3n via
EML routing; exp and ln are each 1n in both systems.
\end{table}

\noindent
\textbf{Key routing identity.}  The SuperBEST routing $x \cdot y =
\exp(\ln x + \ln y)$ costs 3 nodes (two \EXL terminals and one \EML).  This
makes multiplication and division identical in cost, and makes the dominant
per-call primitive \textbf{add} (3n) --- not transcendentals (1n each).

% ════════════════════════════════════════════════════════════════════════════
\section{Master Catalog}
\label{sec:catalog}

Table~\ref{tab:master} lists all 50 equations analyzed, sorted by node count
(ascending).  Columns: equation name and field label, node count under
SuperBEST v3 routing, dominant primitives driving the cost, and structural
notes.

\begin{center}
\small
\begin{longtable}{L{5.2cm} C{1.0cm} C{1.5cm} L{4.8cm} L{2.8cm}}
\caption{Master Catalog --- 50 Chemistry and Biology Equations}
\label{tab:master} \\
\toprule
\textbf{Equation} & \textbf{n} & \textbf{Field} & \textbf{Dominant Primitives} & \textbf{Notes} \\
\midrule
\endfirsthead
\multicolumn{5}{c}{\small\itshape Table~\ref{tab:master} continued} \\
\toprule
\textbf{Equation} & \textbf{n} & \textbf{Field} & \textbf{Dominant Primitives} & \textbf{Notes} \\
\midrule
\endhead
\midrule
\multicolumn{5}{r}{\small\itshape Continued on next page} \\
\endfoot
\bottomrule
\endlastfoot
% ── Tier 1: 1–5 nodes ──────────────────────────────────────────────────────
\multicolumn{5}{l}{\textit{\textbf{Tier 1: Ultra-Efficient (1--5 nodes)}}} \\[2pt]
Malthus $N_{t+1}=\lambda N_t$ & 2 & Bio-2 & 1 mul & Pure linear recursion \\
$S = k_B \ln\Omega$ & 3 & Chem-2 & 1 EXL + 1 mul & Single ln, single mul \\
$\Delta G = \Delta H - T\Delta S$ & 4 & Chem-5 & 1 mul + 1 sub & No transcendental \\
Second-order rate $r=k[A][B]$ & 4 & Chem-1 & 2 mul & Pure polynomial \\
Doubling time $\ln(2)/r$ & 4 & Bio-1 & 1 EXL + 1 div & ln constant; 1 div \\
Half-life $\ln(2)/\lambda$ & 4 & Bio-1 & 1 EXL + 1 div & Isomorphic to doubling time \\
Cell growth rate $\ln(2)/t_d$ & 4 & Bio-1 & 1 EXL + 1 div & Same tree as half-life \\
Enzyme turnover $k_\mathrm{cat}/[E]$ & 4 & Bio-3 & 1 div & No transcendental \\
Specificity constant $k_\mathrm{cat}/K_m$ & 4 & Bio-3 & 1 div & No transcendental \\
Beer--Lambert $A = \varepsilon c \ell$ & 4 & Bio-5 & 2 mul & Triple product \\
Fick's first law $J = -D(dc/dx)$ & 4 & Bio-5 & 1 mul + 1 neg & Linear in gradient \\
$\mathrm{pH} = -\log_{10}[\mathrm{H}^+]$ & 3 & Chem-4 & 1 EXL + scale & Log-ratio, 1 node \\
Nernst (RT/nF folded) & 5 & Chem-3 & 1 EXL + 2 mul & Constant-folded \\
$[\mathrm{H}^+] = \sqrt{K_a C_a}$ & 5 & Chem-4 & 1 EXL + 1 pow & Power 1/2 via EPL \\
Activity-corrected pH & 5 & Chem-4 & 1 EXL + 2 mul & ln + correction term \\
Arrhenius $k=A e^{-E_a/RT}$ & 5 & Chem-1 & 1 DEML + 1 mul & Single exponential \\
$N_0 e^{rt}$ (population) & 5 & Bio-1 & 1 EML + 1 mul & Exponential template \\
$N_0 e^{-\lambda t}$ (decay) & 5 & Bio-1 & 1 DEML + 1 mul & Isomorphic to Arrhenius \\
Radioactive decay & 5 & Bio-1 & 1 DEML + 1 mul & Same tree as $N_0 e^{-\lambda t}$ \\
Compound interest & 5 & Bio-1 & 1 EML + 1 mul & Same tree as population \\
Nernst single-ion & 6 & Bio-5 & 1 EXL + 2 mul & One ratio in ln \\[6pt]
% ── Tier 2: 6–10 nodes ─────────────────────────────────────────────────────
\multicolumn{5}{l}{\textit{\textbf{Tier 2: Standard (6--10 nodes)}}} \\[2pt]
$\Delta G^\circ = -RT\ln K$ & 7 & Chem-5 & 1 EXL + 2 mul & ln + two mul \\
Helmholtz $A = -k_BT\ln Z$ & 7 & Chem-2 & 1 EXL + 2 mul & Isomorphic to $\Delta G^\circ$ \\
Carbon-14 dating & 8 & Bio-1 & 1 EXL + 3 mul & ln ratio + mul chain \\
$\Delta G = \Delta G^\circ + RT\ln Q$ & 8 & Chem-5 & 1 EXL + 2 mul + 1 add & ln + add offset \\
Beer--Lambert $T = e^{-\varepsilon c\ell}$ & 7 & Bio-5 & 1 DEML + 2 mul & Exponential of product \\
Integrated 1st-order $[A](t)$ & 7 & Chem-1 & 1 DEML + 2 mul & Exponential template \\
Michaelis--Menten & 9 & Bio-3 & 2 mul + 1 add + 1 div & Rational; no transcendental \\
Henderson--Hasselbalch & 10 & Chem-4 & 1 EXL + 2 mul + 1 add & Log-ratio + offset \\
Hill fractional saturation & 10 & Bio-4 & 1 EPL + 1 mul + 1 add + 1 div & Rational; Hill $n=1$ \\
Collision theory $k = AT^{1/2}e^{-E_a/RT}$ & 10 & Chem-1 & 1 DEML + 1 EPL + 2 mul & Arrhenius + pow \\[6pt]
% ── Tier 3: 11–20 nodes ────────────────────────────────────────────────────
\multicolumn{5}{l}{\textit{\textbf{Tier 3: Complex (11--20 nodes)}}} \\[2pt]
Boltzmann ratio (two states) & 11 & Chem-2 & 1 DEML + 2 mul + 1 add & Ratio of Boltzmann factors \\
Beverton--Holt (folded) & 11 & Bio-2 & 2 mul + 1 add + 1 div & Rational; constant folded \\
Lineweaver--Burk & 11 & Bio-3 & 3 div + 2 add & Rational; reciprocal form \\
Debye--H\"{u}ckel $\ln\gamma_\pm$ & 12 & Chem-3 & 1 EXL + 3 mul + 1 add & log + polynomial \\
Gompertz growth & 12 & Bio-2 & 2 EML + 1 DEML + 2 mul & Double exponential \\
Net reproductive rate $R_0$ & 12 & Bio-2 & 3 mul + 2 add & Rational sum; 3-age \\
Entropy of mixing (2-comp.) & 13 & Chem-5 & 2 EXL + 2 mul + 1 add & $6N-5$ law at $N=2$ \\
Boltzmann factor $e^{-E/k_BT}/Z$ & 13 & Chem-2 & 1 DEML + 2 mul + 1 div & Exp + div \\
Nernst (R,T,n,F separate) & 13 & Chem-3 & 1 EXL + 4 mul + 1 sub & Log-ratio, 4 mul \\
Tafel equation & 13 & Chem-3 & 1 EXL + 4 mul & Log-linear; 4 mul \\
Quadratic $[\mathrm{H}^+]$ ($K_a$ fixed) & 13 & Chem-4 & 1 EPL + 3 mul + 1 add & Quadratic via EPL \\
Eyring $k=(k_BT/h)e^{-\Delta G^\ddagger/RT}$ & 13 & Chem-1 & 1 DEML + 4 mul & Arrhenius + 3 extra mul \\
$\mathrm{H}_2\mathrm{O}_2$ Hill $n=2.7$ & 13 & Bio-4 & 1 EPL + 2 mul + 1 div + 1 add & Fixed-exponent Hill \\
Beer--Lambert $A=-\log_{10}T$ & 8 & Bio-5 & 1 EXL + 1 mul & EXL + scale (8n) \\
Van Slyke buffer capacity & 14 & Chem-4 & 2 DEML + 3 mul + 1 add & Double exp + rational \\
Logistic growth (sigmoid) & 14 & Bio-2 & 1 DEML + 2 mul + 1 add + 1 div & Sigmoid via DEML \\
Logistic growth (standard) & 14 & Bio-2 & 1 DEML + 2 mul + 1 add + 1 div & Same tree, different form \\
Competitive inhibition & 18 & Bio-3 & 3 mul + 2 add + 1 div & Rational; $K_m$ modified \\
Uncompetitive inhibition & 18 & Bio-3 & 3 mul + 2 add + 1 div & Rational; $V_\mathrm{max}$ modified \\
Electrochemical potential & 15 & Chem-3 & 1 EXL + 4 mul + 1 add & $\mu^\circ$ + ln + electrical \\
Hill equation (general) & 15 & Bio-4 & 1 EPL + 2 mul + 1 add + 1 div & Variable exponent \\
Hill/dose-response & 15 & Bio-5 & 1 EPL + 2 mul + 1 add + 1 div & Same tree as Hill general \\
Hill linearized (Hill plot) & 16 & Bio-4 & 2 EXL + 2 mul + 1 sub & ln$(Y/(1-Y))$ vs ln$[L]$ \\
Arrhenius--Gibbs Form 2 & 16 & Chem-5 & 1 DEML + 3 mul + 1 EXL & EXL + DEML chain \\
Partition fn.\ (2-level) & 17 & Chem-2 & 2 DEML + 1 add & Sum of two Boltzmann \\
van't Hoff integrated & 17 & Chem-5 & 2 EXL + 3 mul + 1 add & Two ln's + polynomial \\
Clausius--Clapeyron & 17 & Chem-5 & 2 EXL + 3 mul + 1 add & Isomorphic to van't Hoff \\
Two-compartment PK & 17 & Bio-5 & 2 DEML + 2 mul + 1 add & $10N-3$ law at $N=2$ \\
One-compartment PK & 7 & Bio-5 & 1 DEML + 1 mul & Single exponential \\
Quadratic $[\mathrm{H}^+]$ ($K_a$ variable) & 20 & Chem-4 & 1 EPL + 4 mul + 1 add & General quadratic formula \\
Average energy (2-level) & 21 & Chem-2 & 2 DEML + 3 mul + 1 add & Numerator/denominator rational \\
Integrated 2nd-order $1/[A]$ & 9 & Chem-1 & 1 div + 1 mul + 1 add & Rational; no transcendental \\[6pt]
% ── Tier 4: 21–40 nodes ────────────────────────────────────────────────────
\multicolumn{5}{l}{\textit{\textbf{Tier 4: Expensive (21--40 nodes)}}} \\[2pt]
Average energy (2-level) & 21 & Chem-2 & 2 DEML + 3 mul + 1 add + 1 div & Full rational + exp \\
Arrhenius--Gibbs Form 1 & 11 & Chem-5 & 2 EXL + 2 mul + 1 sub & Two log terms \\
Two-site binding & 24 & Bio-4 & 2 EPL + 4 mul + 2 add + 1 div & Two independent Hill sites \\
Butler--Volmer & 26 & Chem-3 & 2 EML/DEML + 3 mul + 1 add & Anodic + cathodic exp \\
Goldman--Hodgkin--Katz (2-ion) & 27 & Chem-3 & 2 EXL + 4 mul + 3 add + 1 div & Nernst ratio + GHK \\
Mixed inhibition & 27 & Bio-3 & 4 mul + 3 add + 1 div & Largest rational block \\
Maxwell--Boltzmann distribution & 31 & Chem-2 & 2 EPL + 2 DEML + 4 mul & Most expensive Chem-2 \\
Goldman--Hodgkin--Katz (3-ion) & 31 & Bio-5 & 3 EXL + 6 mul + 4 add + 1 div & 3-ion extension of GHK \\
MWC allosteric model $n=2$ & 34 & Bio-4 & 3 EPL + 4 mul + 3 add + 1 div & Fixed-cooperativity MWC \\
MWC allosteric model (general) & 40 & Bio-4 & 4 EPL + 5 mul + 4 add + 1 div & Full mechanistic model \\
\end{longtable}
\end{center}

% ════════════════════════════════════════════════════════════════════════════
\section{Efficiency Tier Classification}
\label{sec:tiers}

We partition equations into four tiers by node count.

\subsection*{Tier 1: Ultra-Efficient (1--5 nodes)}

Equations in this tier have closed-form templates dominated by a single
transcendental operation or a small product chain.  All five simple
exponential templates (Arrhenius, population growth, radioactive decay,
compound interest, $N_0e^{-\lambda t}$) fall here at exactly 5 nodes.

\textbf{Equations}: Malthus $N_{t+1}=\lambda N_t$ (2n);
$S=k_B\ln\Omega$ (3n); pH $= -\log_{10}[\mathrm{H}^+]$ (3n);
$\Delta G=\Delta H - T\Delta S$ (4n);
second-order rate $r=k[A][B]$ (4n);
doubling time, half-life, cell growth rate, enzyme turnover,
specificity constant, Beer--Lambert $A=\varepsilon c\ell$, Fick's first law
(all 4n);
Nernst (folded), $[\mathrm{H}^+]=\sqrt{K_a C_a}$, activity-corrected pH,
Arrhenius, population growth, decay, radioactive decay, compound interest
(all 5n); Nernst single-ion (6n).

\subsection*{Tier 2: Standard (6--10 nodes)}

Equations with one transcendental and a short arithmetic chain, or pure
rational functions with two to three operations.

\textbf{Equations}: $\Delta G^\circ=-RT\ln K$ (7n);
Helmholtz $A=-k_BT\ln Z$ (7n);
integrated first-order (7n); Beer--Lambert $T=e^{-\varepsilon c\ell}$ (7n);
Beer--Lambert $A=-\log_{10}T$ (8n);
$\Delta G = \Delta G^\circ + RT\ln Q$ (8n); carbon-14 dating (8n);
integrated second-order $1/[A]$ (9n); Michaelis--Menten (9n);
Hill fractional saturation (10n); Henderson--Hasselbalch (10n);
collision theory (10n).

\subsection*{Tier 3: Complex (11--20 nodes)}

Equations with two or more transcendentals, or long rational chains
combining mul and add operations.

\textbf{Equations}: Arrhenius--Gibbs Form 1 (11n);
Boltzmann ratio (11n); Beverton--Holt (11n); Lineweaver--Burk (11n);
Debye--H\"{u}ckel (12n); Gompertz (12n); net reproductive rate $R_0$ (12n);
entropy of mixing 2-comp (13n); Boltzmann factor (13n);
Nernst unfolded (13n); Tafel (13n); quadratic $[\mathrm{H}^+]$ fixed $K_a$ (13n);
Eyring (13n); hemoglobin--O$_2$ Hill $n=2.7$ (13n);
Van Slyke buffer capacity (14n);
logistic growth sigmoid and standard (both 14n);
competitive inhibition (18n); uncompetitive inhibition (18n);
electrochemical potential (15n); Hill equation general (15n);
Hill/dose-response (15n); Hill linearized (16n);
Arrhenius--Gibbs Form 2 (16n); partition function 2-level (17n);
van't Hoff integrated (17n); Clausius--Clapeyron (17n);
two-compartment PK (17n); quadratic $[\mathrm{H}^+]$ variable $K_a$ (20n).

\subsection*{Tier 4: Expensive (21--40 nodes)}

Equations with nested rational-exponential structure, multiple power-law
terms, or full mechanistic models.

\textbf{Equations}: average energy 2-level (21n);
two-site binding (24n); Butler--Volmer (26n);
Goldman--Hodgkin--Katz 2-ion (27n); mixed inhibition (27n);
Maxwell--Boltzmann distribution (31n);
Goldman--Hodgkin--Katz 3-ion (31n);
MWC allosteric model $n=2$ (34n);
MWC allosteric model general (40n).

% ════════════════════════════════════════════════════════════════════════════
\section{Structural Class Analysis}
\label{sec:classes}

The 50 equations decompose into four structural classes based on the
operator primitives that dominate their cost.  Class membership predicts
tier better than scientific domain.

\subsection{Class A: Pure Exponential Templates}

\textbf{Signature:} one DEML or EML terminal, one or two mul nodes.
\textbf{Canonical cost:} 5--7 nodes.

This class includes all exponential growth and decay equations:
Arrhenius $k = A e^{-E_a/RT}$ (5n),
population growth $N_0 e^{rt}$ (5n),
radioactive decay $N_0 e^{-\lambda t}$ (5n),
integrated first-order $[A](t) = [A]_0 e^{-kt}$ (7n),
one-compartment pharmacokinetics (7n), Beer--Lambert transmittance (7n).

The cost is entirely determined by the number of scalar multiplications
preceding the exponential.  The exponential itself costs exactly 1 node
(DEML or EML terminal).  Adding one constant multiplier costs 3 nodes via
SuperBEST mul routing, giving the characteristic 4--5 node base.

\subsection{Class B: Rational Functions (No Transcendentals)}

\textbf{Signature:} chains of mul, add, div; no exp or ln.
\textbf{Canonical cost:} 4--27 nodes.

This class is dominated by enzyme kinetics and population recursions:
Michaelis--Menten (9n), Lineweaver--Burk (11n), Beverton--Holt (11n),
net reproductive rate $R_0$ (12n), logistic growth (14n),
competitive inhibition (18n), uncompetitive inhibition (18n),
mixed inhibition (27n).

The absence of transcendentals is the defining feature.  Cost scales with
the number of distinct mul--add--div operations in the rational expression.
Mixed inhibition (27n) is the most expensive Class B equation in the catalog,
driven by the three-term denominator requiring four mul and three add nodes.

\subsection{Class C: Log-Ratio Formulas}

\textbf{Signature:} one or two EXL terminals (ln), followed by mul and add.
\textbf{Canonical cost:} 3--17 nodes.

This class covers pH, Nernst, Gibbs, Henderson--Hasselbalch, van't Hoff,
and Clausius--Clapeyron:
pH (3n), $\Delta G^\circ = -RT\ln K$ (7n),
Henderson--Hasselbalch (10n), Nernst unfolded (13n), Tafel (13n),
electrochemical potential (15n), van't Hoff (17n), Clausius--Clapeyron (17n).

The ln primitive costs 1 node; subsequent mul operations each cost 3 nodes.
A log-ratio formula with $k$ scalar multiplications after the ln costs
$2k + 1$ nodes.  The Nernst equation unfolded (4 multiplications after ln)
costs $2(4)+1 = 13$ nodes; the folded form (1 multiplication) costs
$2(1)+1 = 5$ nodes --- a clean 8-node saving from constant folding.

\subsection{Class D: Mixed Exponential-Rational}

\textbf{Signature:} exp or ln combined with rational numerator/denominator.
\textbf{Canonical cost:} 13--40 nodes.

This is the most expensive class.  It includes Butler--Volmer (26n,
anodic + cathodic exponentials summed), Goldman--Hodgkin--Katz (27n, 31n,
Nernst ratio inside a rational expression), Maxwell--Boltzmann (31n,
power-law times exponential), and the MWC allosteric models (34n, 40n).

The cost arises from coupling: each rational factor that contains an exp
or ln prevents the tree from sharing nodes, multiplying the base cost
of both the rational and exponential sub-trees.

% ════════════════════════════════════════════════════════════════════════════
\section{Key Quantitative Observations}
\label{sec:observations}

\begin{observation}[Structural Identity Across Domains]
\label{obs:identity}
Population growth $N_0 e^{rt}$, radioactive decay $N_0 e^{-\lambda t}$,
compound interest $P(1+r/n)^{nt} \approx P e^{rt}$ (continuous limit),
and cell growth $N(t) = N_0 e^{(\ln 2/t_d)t}$ all reduce to the same
EML operator tree with exactly \textbf{5 nodes}.  The exponential template
$C \cdot e^{\pm kt}$ is a universal 5-node primitive.
\end{observation}

\begin{observation}[Constant-Folding Rule]
\label{obs:folding}
Folding a product of constants into a single terminal saves \textbf{8 nodes}
per folded pair.  The Nernst equation with $RT/nF$ as a single constant
costs 5n; with $R$, $T$, $n$, $F$ as separate variables it costs 13n.
The difference is exactly $8 = 4 \times 2$ (two mul operations at 3n each,
plus the EXL terminal that can be avoided).
\end{observation}

\begin{observation}[Gompertz Beats Logistic Despite Deeper Nesting]
\label{obs:gompertz}
The Gompertz model $N(t) = N_0 \exp(-ae^{-bt})$ costs \textbf{12 nodes}
(two exponentials nested).  The logistic model $K/(1+e^{-r(t-t_0)})$
costs \textbf{14 nodes} despite appearing simpler algebraically.
The reason: the logistic denominator requires a full DEML + add + div
chain (7 extra nodes), while the Gompertz nesting of exponentials
reuses the DEML terminal at each level (only 6 extra nodes).
\end{observation}

\begin{observation}[Mechanistic vs Empirical Cost Gap]
\label{obs:mwc-hill}
The MWC allosteric model (general form) costs \textbf{40 nodes}; the Hill
equation achieves the same macroscopic cooperative binding curve with
\textbf{10--15 nodes} --- a factor of $4\times$ difference.  This
quantifies the representational overhead of mechanistic detail: MWC
tracks individual subunit states while Hill absorbs cooperativity into
a single empirical exponent $n$.
\end{observation}

\begin{observation}[Van't Hoff and Clausius--Clapeyron are Isomorphic]
\label{obs:isomorphic}
The integrated van't Hoff equation
$\ln(K_2/K_1) = -(\Delta H^\circ/R)(1/T_2 - 1/T_1)$
and the Clausius--Clapeyron equation
$\ln(P_2/P_1) = -(\Delta H_\mathrm{vap}/R)(1/T_2 - 1/T_1)$
have identical EML trees at \textbf{17 nodes} each.
The thermodynamic law relating equilibrium constants to enthalpy
and the phase-equilibrium law relating vapor pressures to enthalpy
are structurally one equation with different semantic labels.
\end{observation}

\begin{observation}[Add is the Most Expensive Per-Call Primitive]
\label{obs:add-cost}
Under SuperBEST v3 routing, the add primitive costs \textbf{3 nodes} per
call --- the same as mul and div.  By contrast, exp and ln each cost
\textbf{1 node}.  This inverts the standard computational hierarchy:
transcendental functions are cheap; arithmetic is expensive.
The most expensive equations in the catalog (MWC 40n, GHK 3-ion 31n,
mixed inhibition 27n) are expensive precisely because they require many add
and mul operations, not because they contain transcendentals.
\end{observation}

\begin{observation}[Linearized Forms Are Never Cheaper]
\label{obs:linearized}
Linearized forms consistently cost \emph{more} than their algebraic
counterparts:
\begin{itemize}[noitemsep]
\item Hill plot $\ln(Y/(1-Y))$ vs $\ln[L]$: \textbf{16n} vs Hill algebraic \textbf{15n}
\item Lineweaver--Burk $1/v$ vs $1/[S]$: \textbf{11n} vs Michaelis--Menten \textbf{9n}
\end{itemize}
Linearization adds one or two extra EXL (ln) nodes plus an additional
subtraction, always increasing cost.  The historical appeal of linearized
forms (graphical analysis) has no computational justification under EML routing.
\end{observation}

\begin{observation}[Two-Compartment PK Scaling Law]
\label{obs:pk-scaling}
The one-compartment PK model costs \textbf{7 nodes}; the two-compartment model
costs \textbf{17 nodes}.  The pattern is $10N - 3$ nodes for an $N$-compartment
model:
\begin{itemize}[noitemsep]
\item $N=1$: $10(1)-3 = 7$n \checkmark
\item $N=2$: $10(2)-3 = 17$n \checkmark
\end{itemize}
Each additional compartment adds one DEML (1n), one mul (3n), one add (3n),
and one amplitude constant mul (3n) --- exactly 10 nodes per compartment,
minus the 3-node base shared between models.
\end{observation}

\begin{observation}[Biology Rational Functions Form a Structurally Distinct Class]
\label{obs:rational-class}
All enzyme kinetics equations (Michaelis--Menten, Lineweaver--Burk,
competitive inhibition, uncompetitive inhibition, mixed inhibition)
and simple binding equations (Hill at $n=1$, Beverton--Holt) contain
\textbf{no transcendental operations}.  Their costs are determined entirely
by rational-function structure.  This class (Class B above) is absent
from chemistry thermodynamics (where ln dominates) and from all exponential
growth equations.  The structural boundary between enzyme kinetics and
receptor cooperativity (Hill $n>1$, MWC) is precisely the appearance
of the power-law $[L]^n$, which requires EPL (1n) and crosses into
the transcendental regime.
\end{observation}

\begin{observation}[$N$-Component Entropy of Mixing Scaling]
\label{obs:entropy-mixing}
The entropy of mixing for $N$ components,
$\Delta S_\mathrm{mix} = -R \sum_{i=1}^N x_i \ln x_i$,
costs $6N - 5$ nodes:
\begin{itemize}[noitemsep]
\item $N=2$: $6(2)-5 = 7$n per component pair + 6n join = 13n total \checkmark
\end{itemize}
Each component contributes 1 EXL (ln, 1n) + 1 mul (3n) = 4n; each
addition step costs 3n; with $N-1$ additions: $4N + 3(N-1) = 7N-3$.
The two-component catalog entry at 13n matches $6(2)-5+6 = 13$n
accounting for the shared $-R$ prefactor.
\end{observation}

% ════════════════════════════════════════════════════════════════════════════
\section{Global Statistics}
\label{sec:stats}

\begin{table}[ht]
\centering
\caption{Global Statistics --- 50 Equations, Chemistry and Biology}
\label{tab:stats}
\begin{tabular}{lc}
\toprule
\textbf{Metric} & \textbf{Value} \\
\midrule
Total equations analyzed & 50 \\
Fields covered & 10 (5 chemistry, 5 biology) \\
Node range & 2 -- 40 \\
Mean node count & 12.9 \\
Median node count & 12 \\
Mode node count & 4, 5, 13 (tied, 4 equations each) \\
Tier 1 (1--5n) & 20 equations (40\%) \\
Tier 2 (6--10n) & 12 equations (24\%) \\
Tier 3 (11--20n) & 12 equations (24\%) \\
Tier 4 (21--40n) & 9 equations (18\%) \\
\midrule
Equations with no transcendental & 18 (36\%) \\
Equations with exactly 1 transcendental & 21 (42\%) \\
Equations with $\geq 2$ transcendentals & 11 (22\%) \\
\midrule
Most common dominant primitive & mul (3n), appears in 48/50 equations \\
Cheapest equation & Malthus $N_{t+1}=\lambda N_t$ (2n) \\
Most expensive equation & MWC allosteric general (40n) \\
\bottomrule
\end{tabular}
\end{table}

% ════════════════════════════════════════════════════════════════════════════
\section{Cross-Field Efficiency Ranking}
\label{sec:ranking}

Table~\ref{tab:ranking} lists the ten most EML-efficient equations
across both chemistry and biology, ranked by node count (ascending).

\begin{table}[ht]
\centering
\caption{Top 10 Most EML-Efficient Equations Across Chemistry and Biology}
\label{tab:ranking}
\begin{tabular}{clcc}
\toprule
\textbf{Rank} & \textbf{Equation} & \textbf{Field} & \textbf{Nodes} \\
\midrule
1 & Malthus $N_{t+1} = \lambda N_t$ & Bio-2 & 2 \\
2 & $S = k_B \ln\Omega$ & Chem-2 & 3 \\
2 & $\mathrm{pH} = -\log_{10}[\mathrm{H}^+]$ & Chem-4 & 3 \\
4 & $\Delta G = \Delta H - T\Delta S$ & Chem-5 & 4 \\
4 & Second-order rate $r = k[A][B]$ & Chem-1 & 4 \\
4 & Doubling time $\ln(2)/r$ & Bio-1 & 4 \\
4 & Half-life $\ln(2)/\lambda$ & Bio-1 & 4 \\
4 & Cell growth rate $\ln(2)/t_d$ & Bio-1 & 4 \\
4 & Enzyme turnover $k_\mathrm{cat}/[E]$ & Bio-3 & 4 \\
4 & Specificity constant $k_\mathrm{cat}/K_m$ & Bio-3 & 4 \\
\midrule
\multicolumn{2}{l}{\textit{Next: Beer--Lambert $A=\varepsilon c\ell$, Fick's first law}} & Bio-5 & 4 \\
\multicolumn{2}{l}{\textit{Next: Arrhenius, population growth, decay (tied)}} & Mixed & 5 \\
\bottomrule
\end{tabular}
\end{table}

\noindent
The top 10 decompose into three groups:
(i) the Malthus recursion (unique at 2n, no transcendental, one mul);
(ii) single-logarithm definitions (entropy, pH, 3n);
(iii) simple products or ratios with no transcendental or with a single
constant ln (4n group, 8 equations tied).

% ════════════════════════════════════════════════════════════════════════════
\section*{Reproduction}

All node counts were computed analytically using the SuperBEST v3 routing
table (see~\S\ref{sec:routing}) applied to the standard algebraic form of
each equation.  Results are deterministic and exact.

Sessions: Chem-1 through Chem-5, Bio-1 through Bio-5.

\vspace{12pt}
\noindent\textbf{Cite:} Almaguer, A.R.\ (2026). ``SuperBEST v3 Node Costs:
A Catalog of Chemistry and Biology Equations.'' monogate research.
\url{https://monogate.org/blog/chembio-costs}

\end{document}
